\paragrafos{Κυκλικά διαγράμματα και ιστογράμματα}{Γ.1.3, «Κυκλικά διαγράμματα – Ιστογράμματα συχνοτήτων κλάσεων ίσου πλάτους με δεδομένο το πλήθος των κλάσεων»}

\begin{anagke}
Η θεία Μαρία έφτιαξε μια πίτσα για τα τρία ανίψια της που είναι $4$, $8$ και $12$ ετών. Θέλει να μοιράσει την πίτσα ανάλογα με την ηλικία τους. Όλοι κοιτούν την πίτσα: τα $4$, $8$ και $12$ κάνουν $24$ μαζί, άρα η μικρότερη πρέπει να πάρει $\frac{4}{24}$ της πίτσας. Αλλά πώς κόβεις ένα κομμάτι που είναι $\frac{4}{24}$ από μια στρογγυλή πίτσα; «Με γωνία από το κέντρο», λέει ο Ηλίας, «το $\frac{4}{24}$ του κύκλου».

Την ίδια ώρα η Ζωή έχει ένα άλλο πρόβλημα. Η παιδίατρος της γειτονιάς σημείωσε την ηλικία των παιδιών που πέρασαν από το ιατρείο μια μέρα και ο πίνακας έχει πάρα πολλές τιμές. «Θέλω ένα σχήμα που να φαίνεται πού μαζεύονται οι ηλικίες», λέει η Ζωή.

Σε αυτή την παράγραφο μαθαίνουμε δύο νέους τρόπους να δείχνουμε δεδομένα με εικόνα: το κυκλικό διάγραμμα, που μοιάζει με πίτσα, και το ιστόγραμμα.
\end{anagke}

\begin{lexi}{κυκλικό, τομέας, επίκεντρη, ιστόγραμμα}{ονόματα για τις εικόνες}
\lexiword{Κυκλικό} είναι ό,τι έχει σχήμα κύκλου. \lexiword{Τομέας} είναι το κομμάτι της πίτσας, από το κέντρο ως την άκρη. Η γωνία του τομέα λέγεται \lexiword{επίκεντρη}, γιατί η κορυφή της είναι στο κέντρο του κύκλου. Το \lexiword{ιστόγραμμα} παίρνει το όνομά του από τον \emph{ιστό}, που σημαίνει «όρθιο στύλο» (όπως ο ιστός ενός πλοίου): τα ορθογώνια του διαγράμματος στέκονται όρθια, σαν στύλοι, το ένα δίπλα στο άλλο.
\end{lexi}

\begin{metafrasi}
Ένα κομμάτι πίτσας. & τομέας κύκλου\\
Η γωνία του κομματιού στο κέντρο. & επίκεντρη γωνία $\hat{\omega}$\\
Ολόκληρος ο κύκλος. & $360^\circ$ και $100\%$\\
Μισή πίτσα. & $180^\circ$ και $50\%$\\
Ορθογώνια κολλητά, ένα για κάθε κλάση. & ιστόγραμμα\\
Πόσο ψηλό είναι το ορθογώνιο. & η συχνότητα της κλάσης\\
\end{metafrasi}

\tmima{Η θεωρία}

\ypotmima{Α}{Κυκλικό διάγραμμα}

\begin{orismos}{κυκλικό διάγραμμα}
Το \textbf{κυκλικό διάγραμμα} είναι ένας κυκλικός δίσκος χωρισμένος σε κυκλικούς τομείς, που παριστάνουν τις διαφορετικές τιμές μιας μεταβλητής.
\end{orismos}

Οι επίκεντρες γωνίες των τομέων είναι ανάλογες με τις συχνότητες των τιμών. Ολόκληρος ο κύκλος αντιστοιχεί στο σύνολο των παρατηρήσεων και έχει γωνία $360^\circ$.

\begin{kanonas}{επίκεντρη γωνία}
$$\frac{\text{συχνότητα της τιμής}}{\text{σύνολο}}=\frac{\hat{\omega}}{360^\circ}\qquad\text{άρα}\qquad \hat{\omega}=\frac{\text{συχνότητα}}{\text{σύνολο}}\cdot 360^\circ$$
Με τη σχετική συχνότητα: $\hat{\omega}=\text{σχετική συχνότητα}\cdot360^\circ$. Το άθροισμα όλων των επίκεντρων γωνιών είναι $360^\circ$.
\end{kanonas}

Το κυκλικό διάγραμμα είναι χρήσιμο για ποσοστά και για λίγες κατηγορίες. Έχει όμως και μειονεκτήματα: είναι δύσκολο να συγκρίνουμε δύο τομείς που διαφέρουν λίγο και δύσκολο να διακρίνουμε πολύ μικρούς τομείς. Όταν οι τιμές είναι πολλές, δεν ταιριάζει.

\ypotmima{Β}{Ιστόγραμμα συχνοτήτων}

Όταν τα δεδομένα έχουν ομαδοποιηθεί σε κλάσεις ίσου πλάτους, τα παριστάνουμε με \textbf{ιστόγραμμα}. 

\begin{orismos}{ιστόγραμμα}
Το \textbf{ιστόγραμμα συχνοτήτων} αποτελείται από συνεχόμενα ορθογώνια που έχουν \textbf{βάση} ίση με το πλάτος κάθε κλάσης και \textbf{ύψος} ίσο με τη συχνότητα της αντίστοιχης κλάσης. Τα άκρα των κλάσεων γράφονται στον οριζόντιο άξονα. Με τον ίδιο τρόπο φτιάχνεται και το \textbf{ιστόγραμμα σχετικών συχνοτήτων}, με ύψη τις σχετικές συχνότητες.
\end{orismos}

Το ιστόγραμμα ταιριάζει σε μεγάλο αριθμό συνεχών δεδομένων που έχουν ομαδοποιηθεί: τα ορθογώνια είναι κολλητά, γιατί το τέλος μιας κλάσης είναι η αρχή της επόμενης, και μας δείχνουν με μια ματιά πού μαζεύονται οι τιμές. Δεν ταιριάζει για κατηγορικά δεδομένα (εκεί χρησιμοποιούμε ραβδόγραμμα) και η εικόνα αλλάζει ανάλογα με τον αριθμό των κλάσεων που θα διαλέξουμε.

\begin{prosoxi}
\begin{itemize}[leftmargin=1.2em,itemsep=0.3ex,topsep=2pt]
\item Στο ραβδόγραμμα οι ράβδοι έχουν κενά, στο ιστόγραμμα τα ορθογώνια ακουμπούν.
\item Στο κυκλικό διάγραμμα οι γωνίες βρίσκονται με τον τύπο $\hat{\omega}=\frac{\text{συχνότητα}}{\text{σύνολο}}\cdot360^\circ$. Μην ξεχνάς το $360^\circ$ ή να διαιρέσεις με το σύνολο.
\item Το ύψος του ορθογωνίου στο ιστόγραμμα είναι η συχνότητα, όχι το πλάτος της κλάσης.
\item Το άθροισμα των γωνιών ενός κυκλικού διαγράμματος πρέπει να βγαίνει $360^\circ$. Αν δεν βγαίνει, έχει γίνει λάθος.
\end{itemize}
\end{prosoxi}

\tmima{Λυμένα παραδείγματα}

\paradeigma{Κατασκευή κυκλικού διαγράμματος}
\tixero Η Φανή ρώτησε $30$ παιδιά τι πίνουν στο διάλειμμα: νερό $12$, χυμό $9$, γάλα $6$, κάτι άλλο $3$.\par
\tithelo Την επίκεντρη γωνία κάθε τομέα και το κυκλικό διάγραμμα.\par
\possindeo Με τον τύπο $\hat{\omega}=\frac{\text{συχνότητα}}{30}\cdot360^\circ$: νερό $\frac{12}{30}\cdot360^\circ=144^\circ$, χυμός $\frac{9}{30}\cdot360^\circ=108^\circ$, γάλα $\frac{6}{30}\cdot360^\circ=72^\circ$, άλλο $\frac{3}{30}\cdot360^\circ=36^\circ$.\par
\apantisi

\begin{center}
\renewcommand{\arraystretch}{1.25}
\begin{tabular}{lccc}
\toprule
\textbf{Ρόφημα} & \textbf{Συχνότητα} & \textbf{Σχετική συχνότητα} & \textbf{Γωνία}\\
\midrule
Νερό & $12$ & $40\%$ & $144^\circ$\\
Χυμός & $9$ & $30\%$ & $108^\circ$\\
Γάλα & $6$ & $20\%$ & $72^\circ$\\
Άλλο & $3$ & $10\%$ & $36^\circ$\\
\midrule
\textbf{Σύνολο} & $30$ & $100\%$ & $360^\circ$\\
\bottomrule
\end{tabular}
\end{center}

\begin{center}
\gaKykliko{1.4}{Νερό/12/$144^\circ$,Χυμός/9/$108^\circ$,Γάλα/6/$72^\circ$,Άλλο/3/$36^\circ$}
\end{center}

\elegxos $144+108+72+36=360$ και $40+30+20+10=100$.

\paradeigma{Διαβάζω κυκλικό διάγραμμα}
\tixero Στο παρακάτω διάγραμμα φαίνονται οι $240$ επισκέπτες μιας σχολικής γιορτής. Οι γωνίες των τομέων είναι: ενήλικες $180^\circ$, παιδιά $90^\circ$, ηλικιωμένοι $60^\circ$, άλλοι $30^\circ$.\par
\tithelo Πόσοι επισκέπτες είναι σε κάθε κατηγορία;\par
\possindeo Από τον τύπο, $\text{συχνότητα}=\frac{\hat{\omega}}{360^\circ}\cdot240$. Ενήλικες: $\frac{180}{360}\cdot240=120$. Παιδιά: $\frac{90}{360}\cdot240=60$. Ηλικιωμένοι: $\frac{60}{360}\cdot240=40$. Άλλοι: $\frac{30}{360}\cdot240=20$.\par
\apantisi $120$, $60$, $40$ και $20$ επισκέπτες.

\begin{center}
\gaKykliko{1.4}{Ενήλικες/180/$180^\circ$,Παιδιά/90/$90^\circ$,Ηλικιωμένοι/60/$60^\circ$,Άλλοι/30/$30^\circ$}
\end{center}

\elegxos $120+60+40+20=240$.

\paradeigma{Κατασκευή ιστογράμματος}
\tixero Ο πίνακας των υψών της Α$3$ από την προηγούμενη παράγραφο: οι κλάσεις $145-150$, $150-155$, $155-160$, $160-165$, $165-170$, $170-175$ έχουν συχνότητες $1$, $6$, $6$, $6$, $4$, $2$.\par
\tithelo Το ιστόγραμμα συχνοτήτων και το ιστόγραμμα σχετικών συχνοτήτων.\par
\possindeo Στον οριζόντιο άξονα βάζουμε τα άκρα των κλάσεων. Φτιάχνουμε ένα ορθογώνιο πάνω από κάθε κλάση, με ύψος ίσο με τη συχνότητά της. Για το ιστόγραμμα σχετικών συχνοτήτων βάζουμε ύψη $4\%$, $24\%$, $24\%$, $24\%$, $16\%$, $8\%$.\par
\apantisi

\begin{center}
\gaIstogramma{145}{5}{6}{1}{1,6,6,6,4,2}{ύψος (cm)}{συχνότητα}\hfill
\gaIstogramma{145}{5}{24}{4}{4,24,24,24,16,8}{ύψος (cm)}{σχετική συχνότητα (\%)}
\end{center}

\elegxos Οι δύο εικόνες έχουν το ίδιο σχήμα. Αλλάζει μόνο η κλίμακα στον κατακόρυφο άξονα.

\paradeigma{Διαβάζω ιστόγραμμα}
\tixero Το ιστόγραμμα δείχνει πόσα λεπτά περίμεναν στον φούρνο $40$ πελάτες.\par
\begin{center}
\gaIstogramma{0}{2}{14}{2}{6,14,12,6,2}{χρόνος αναμονής (λεπτά)}{πελάτες}
\end{center}
\tithelo Πόσοι περίμεναν λιγότερο από $4$ λεπτά; Ποια κλάση έχει τους περισσότερους; Ποιο ποσοστό περίμενε τουλάχιστον $6$ λεπτά;\par
\possindeo Οι κλάσεις έχουν πλάτος $2$. Λιγότερο από $4$ λεπτά σημαίνει οι κλάσεις $0-2$ και $2-4$: $6+14=20$ πελάτες. Ο ψηλότερος στύλος είναι ο $2-4$, με συχνότητα $14$. Τουλάχιστον $6$ λεπτά σημαίνει οι $6-8$ και $8-10$: $6+2=8$ πελάτες, δηλαδή $\frac{8}{40}=20\%$.\par
\apantisi $20$ πελάτες, η κλάση $2-4$, και $20\%$.\par
\elegxos $6+14+12+6+2=40$.

\tmima{Ασκήσεις}

\zoni{Το κατάλαβα}{διαβάζω και υπολογίζω}

\ex Στο κυκλικό διάγραμμα φαίνεται πώς περνούν το Σάββατο τα $40$ παιδιά μιας τάξης.
\begin{center}
\gaKykliko{1.4}{Παρέα/15/15 παιδιά,Οθόνη/10/10 παιδιά,Σπορ/10/10 παιδιά,Βόλτα/5/5 παιδιά}
\end{center}
α) Ποιο ποσοστό των παιδιών κάνει σπορ; β) Ποια είναι η επίκεντρη γωνία του τομέα «Παρέα»; γ) Ποια η γωνία του τομέα «Βόλτα»;

\ex Ρωτήσαμε $40$ μαθητές πώς πηγαίνουν στο σχολείο: με τα πόδια $14$, με λεωφορείο $10$, με αυτοκίνητο $12$, με ποδήλατο $4$. Βρες τις σχετικές συχνότητες και τις επίκεντρες γωνίες των τομέων.

\ex Σε ένα κυκλικό διάγραμμα για τις δραστηριότητες μιας ομάδας παιδιών οι γωνίες είναι: ζωγραφική $90^\circ$, μουσική $120^\circ$, θέατρο $60^\circ$, χορός $60^\circ$, άλλο $30^\circ$. Στη ζωγραφική συμμετέχουν $12$ παιδιά. Πόσα παιδιά έχει η ομάδα και πόσα συμμετέχουν σε κάθε δραστηριότητα;

\ex Στο ιστόγραμμα φαίνονται οι χρόνοι σε δευτερόλεπτα $40$ παιδιών στα $50$ μέτρα ελεύθερο.
\begin{center}
\gaIstogramma{30}{5}{16}{4}{4,10,16,6,4}{χρόνος (s)}{παιδιά}
\end{center}
α) Ποιο είναι το πλάτος κάθε κλάσης; β) Ποια κλάση έχει τη μεγαλύτερη συχνότητα; γ) Πόσα παιδιά έκαναν λιγότερο από $40$ s; δ) Ποιο ποσοστό έκανε τουλάχιστον $45$ s;

\ex Συμπλήρωσε. α) Στο ιστόγραμμα η βάση κάθε ορθογωνίου ισούται με το \ldots\ της κλάσης και το ύψος με τη \ldots\ της κλάσης. β) Ένας τομέας με σχετική συχνότητα $25\%$ έχει επίκεντρη γωνία \ldots. γ) Το άθροισμα όλων των επίκεντρων γωνιών ενός κυκλικού διαγράμματος είναι \ldots. δ) Τα ορθογώνια του ιστογράμματος \ldots\ μεταξύ τους, ενώ οι ράβδοι του ραβδογράμματος όχι.

\ex Ποια αναπαράσταση ταιριάζει καλύτερα: ραβδόγραμμα, κυκλικό διάγραμμα ή ιστόγραμμα; α) Τα ποσοστά που πήρε κάθε υποψήφιος στις εκλογές του μαθητικού συμβουλίου, με τρεις υποψηφίους. β) Τα ύψη $200$ μαθητών, ομαδοποιημένα σε κλάσεις. γ) Πόσα παιδιά προτιμούν κάθε ένα από $8$ διαφορετικά αθλήματα.

\ex Σωστό ή λάθος; α) Ένα ποσοστό $50\%$ αντιστοιχεί σε τομέα $180^\circ$. β) Σε ιστόγραμμα το ύψος κάθε ορθογωνίου είναι ίσο με το πλάτος της κλάσης. γ) Σε κυκλικό διάγραμμα, για σχετική συχνότητα $25\%$ η επίκεντρη γωνία είναι $90^\circ$. δ) Το κυκλικό διάγραμμα ταιριάζει όταν οι κατηγορίες είναι πάρα πολλές. ε) Ένα ιστόγραμμα συχνοτήτων και ένα ιστόγραμμα σχετικών συχνοτήτων των ίδιων δεδομένων έχουν το ίδιο σχήμα.

\zoni{Το χρησιμοποιώ}{φτιάχνω διαγράμματα}

\ex \textbf{Το φύλλο.} Η θεία Μαρία μέτρησε το μήκος $20$ φύλλων από το δέντρο της. Οι κλάσεις $4-6$, $6-8$, $8-10$, $10-12$, $12-14$ (σε cm) έχουν συχνότητες $2$, $5$, $8$, $4$, $1$. Φτιάξε το ιστόγραμμα συχνοτήτων.

\ex \textbf{Το βιβλιοπωλείο.} Σε $60$ αγορές ενός βιβλιοπωλείου: τετράδια $24$, στυλό $15$, μπλοκ $12$, άλλα είδη $9$. Βρες τις επίκεντρες γωνίες και φτιάξε το κυκλικό διάγραμμα.

\ex \textbf{Οι τσάντες.} Στην παράγραφο Γ1.2 τα $40$ παιδιά είχαν συχνότητες τσάντας $6$, $12$, $14$, $6$, $2$ στις κλάσεις $2-4$, $4-6$, $6-8$, $8-10$, $10-12$. Βρες τις σχετικές συχνότητες και φτιάξε το ιστόγραμμα σχετικών συχνοτήτων.

\ex \textbf{Οι προτιμήσεις.} Στο κυκλικό διάγραμμα φαίνονται οι προτιμήσεις $36$ παιδιών για τρία αθλήματα και μία κατηγορία «άλλο».
\begin{center}
\gaKykliko{1.4}{Ποδόσφαιρο/150/$150^\circ$,Μπάσκετ/90/$90^\circ$,Κολύμβηση/70/$70^\circ$,Άλλο/50/$50^\circ$}
\end{center}
Πόσα παιδιά προτιμούν κάθε άθλημα; Ποιο κλάσμα των παιδιών προτιμά το ποδόσφαιρο;

\ex \textbf{Το μουσείο.} Το ιστόγραμμα δείχνει τις ηλικίες $60$ επισκεπτών ενός μουσείου.
\begin{center}
\gaIstogramma{0}{10}{16}{4}{6,15,12,9,12,6}{ηλικία (έτη)}{επισκέπτες}
\end{center}
α) Πόσοι επισκέπτες είναι κάτω από $20$ ετών και ποιο ποσοστό είναι αυτό; β) Ποιο ποσοστό είναι από $40$ ετών και πάνω; γ) Σε ποια κλάση ανήκει η πιο συνηθισμένη ηλικία;

\zoni{Το σκέφτομαι}{γιατί και πότε}

\ex Η Αγγελική έφτιαξε κυκλικό διάγραμμα για το αγαπημένο χρώμα των συμμαθητών της, με $9$ διαφορετικά χρώματα και μερικά από αυτά με πολύ μικρούς τομείς. Τι δυσκολίες θα έχει όποιος το διαβάσει; Ποιο άλλο διάγραμμα θα πρότεινες;

\ex Εξήγησε γιατί το $1\%$ αντιστοιχεί σε $3{,}6^\circ$ στο κυκλικό διάγραμμα. Βρες την επίκεντρη γωνία για ποσοστά $5\%$, $12\%$ και $37{,}5\%$.

\ex Στο ιστόγραμμα της Α$3$ για τα ύψη, ο Ηλίας ενώνει ανά δύο τις γειτονικές κλάσεις και φτιάχνει τρεις κλάσεις πλάτους $10$: $145-155$, $155-165$, $165-175$. Βρες τις νέες συχνότητες και σχεδίασε το νέο ιστόγραμμα. Τι αλλάζει στην εικόνα και τι μένει ίδιο;

\ex Ο Βασίλης υπολόγισε τις γωνίες για $24$ παιδιά: $8$ παιδιά $120^\circ$, $6$ παιδιά $90^\circ$, $7$ παιδιά $110^\circ$, $3$ παιδιά $45^\circ$. Βρες ποια γωνία είναι λάθος και διόρθωσέ την.

\zoni{Για το σπίτι}{η μέρα μου σε κύκλο}

\ex Χώρισε σε κυκλικό διάγραμμα τις $24$ ώρες της μέρας σου: ύπνος, σχολείο, φαγητό και μετακινήσεις, διάβασμα, ελεύθερος χρόνος. Βρες πρώτα πόσες μοίρες αντιστοιχούν σε μία ώρα.

\ex Βρες σε εφημερίδα ή ιστοσελίδα ένα κυκλικό διάγραμμα ή ένα ιστόγραμμα και γράψε τρία συμπεράσματα που προκύπτουν από αυτό.
